% Supplement: full proofs. Compile: pdflatex supplement (twice).
\documentclass[10pt]{article}
\usepackage[margin=1in]{geometry}
\usepackage{amsmath,amssymb,amsthm,bm}
\newtheorem{theorem}{Theorem}
\newtheorem{lemma}{Lemma}
\newtheorem{proposition}{Proposition}
\newtheorem{remark}{Remark}
\DeclareMathOperator{\tr}{tr}
\DeclareMathOperator{\range}{range}
\newcommand{\R}{\mathbb{R}}
\newcommand{\C}{\mathbb{C}}
\newcommand{\Z}{\mathbb{Z}}
\newcommand{\Lam}{\Lambda}
\newcommand{\Phimat}{\bm{\Phi}}
\newcommand{\norm}[1]{\lVert#1\rVert}
\newcommand{\vt}{v_T}
\newcommand{\at}{a_T}

\title{Supplement to ``Silent Aliasing in Fixed Fourier-Feature Coordinate Models'':\\
Full Proofs}
\author{Weiping Yan}
\date{}

\begin{document}
\maketitle

Numbering follows the main paper (T1--T4 are Theorems 1--4 there).  Setting and notation
as in the main paper: $e_\omega(x)=e^{i2\pi\omega x}$ on $[0,1)$, model set
$\Lam=\{\omega_1,\dots,\omega_m\}$, samples $T=\{t_j\}_{j=1}^N$, synthesis matrix
$\Phimat_{jk}=e_{\omega_k}(t_j)$, sample atom $\bm\phi_\nu=(e^{i2\pi\nu t_j})_j$,
visibility $\vt(\nu)=\norm{(\bm I-\bm P_\Lam)\bm\phi_\nu}/\sqrt N$, aliasability
$\at(\nu)=\norm{\Phimat^\dagger\bm\phi_\nu}$.  Every constant below is also asserted
numerically in \texttt{tests/test\_identifiability.py}.

\section{Proof of Theorem 1 (exact equivalence on grid subsets)}

\textbf{(i)}  $\vt(\nu)=0$ iff $(\bm I-\bm P_\Lam)\bm\phi_\nu=0$ iff
$\bm\phi_\nu\in\range\Phimat$, which is precisely the statement that some coefficient
vector $\bm c$ has $\Phimat\bm c=\bm\phi_\nu$, i.e.\ the unit tone at $\nu$ agrees with a
model element at every sample.  For the grid statement, let $t_j=g_j/Q$ with
$g_j\in\{0,\dots,Q-1\}$ and $\nu=\omega_k+rQ$, $r\in\Z$.  Then
\[
e^{i2\pi\nu g_j/Q}=e^{i2\pi\omega_k g_j/Q}\,e^{i2\pi r g_j}=e^{i2\pi\omega_k g_j/Q},
\]
since $rg_j\in\Z$; hence $\bm\phi_\nu=\bm\phi_{\omega_k}$ entrywise, for \emph{every}
subset of the grid.  Consequently $\vt(\nu)=0$; moreover, because
$\sigma_{\min}(\Phimat)>0$ makes the LS solution unique, the fit of $y=\bm\phi_\nu$ is
the unique $\bm c$ with $\Phimat\bm c=\bm\phi_{\omega_k}$, namely the $k$-th unit vector,
and the residual is exactly zero.

\textbf{(ii)}  By (i), $S_Tf_1=S_Tf_2$ pointwise.  For any noise mechanism that acts on
$S_Tf$ (additive or otherwise, as long as its law depends on the signal only through
$S_Tf$), the observation distributions under $f_1$ and $f_2$ are therefore identical.
Let $\hat f=\delta(\bm y)$ be any estimator.  For every realization $\bm y$,
\[
\norm{\delta(\bm y)-f_1}_{L^2}+\norm{\delta(\bm y)-f_2}_{L^2}
\ \ge\ \norm{f_1-f_2}_{L^2}=|a|\,\norm{e_\nu-e_{\omega_k}}_{L^2}=|a|\sqrt2,
\]
using the triangle inequality and, for the last equality, that $e_\nu$ and
$e_{\omega_k}$ are distinct integer frequencies, hence orthonormal in $L^2[0,1)$.  Thus
$\max_i\norm{\delta(\bm y)-f_i}\ge|a|/\sqrt2$ pointwise (a two-point bound in the sense
of Le Cam, here even realization-wise).

\textbf{(iii)}  If $[\Phimat\;\bm\phi_\nu]$ has full column rank then
$\bm\phi_\nu\notin\range\Phimat$, so $\vt(\nu)>0$ by (i). \hfill$\square$

\begin{remark}
The full-column-rank hypothesis on $\Phimat$ requires $N\ge m$ and the elements of
$\Lam$ to be distinct modulo $Q$ (two atoms congruent mod $Q$ have identical columns on
any grid subset).  ``Any $N$'' claims are always to be read as ``any grid-subset size
$N$ for which this rank condition holds''.
\end{remark}

\section{Proof of Theorem 2 (Riesz conversion and decomposition)}

\textbf{(i)}  $G_{L^2}$ is the Gram matrix of $\{e_{\omega_k}\}$ in $L^2[0,1)$:
$\norm{\sum_kc_ke_{\omega_k}}_{L^2}^2=\bm c^{*}G_{L^2}\bm c$, and the Rayleigh quotient
of the Hermitian PSD matrix $G_{L^2}$ lies in
$[\lambda_{\min},\lambda_{\max}]$.  The entries are
$\int_0^1e^{i2\pi(\omega_\ell-\omega_k)t}\,dt
=\frac{e^{i2\pi\delta}-1}{i2\pi\delta}=e^{i\pi\delta}\,\frac{\sin\pi\delta}{\pi\delta}$
for $\delta=\omega_\ell-\omega_k\neq0$, and $1$ on the diagonal.  $G_{L^2}$ is singular
iff the atoms are linearly dependent in $L^2$, which for finitely many distinct
frequencies never happens (distinct exponentials are linearly independent); hence
$\lambda_{\min}>0$ for distinct frequencies.

\textbf{(ii)}  Expand
$\hat f-f=\sum_k\Delta c_k e_{\omega_k}-\sum_l a_le_{\nu_l}$ and evaluate
$\norm{\hat f-f}_{L^2}^2$ as the quadratic form of the joint Gram matrix of
$(\{\omega_k\},\{\nu_l\})$ applied to $(\Delta\bm c,-\bm a)$; grouping blocks gives
Eq.~(3) of the main paper.  For all-integer frequencies the off-diagonal blocks vanish
(orthonormality), killing the cross term. \hfill$\square$

\section{Proof of Theorem 3 (random sampling breaks exact aliasing)}

We use the following complex Hoeffding bound throughout: if $Z_1,\dots,Z_N$ are
independent with $|Z_j|\le1$, then for $\bar Z=\frac1N\sum_jZ_j$,
\begin{equation}\label{eq:hoeffding}
\Pr\big(|\bar Z-\mathbb E\bar Z|\ge\varepsilon\big)\ \le\
4\exp(-N\varepsilon^2/4),
\end{equation}
obtained by applying the real Hoeffding inequality to the real and imaginary parts
(each in $[-1,1]$) with radius $\varepsilon/\sqrt2$ and a union bound.

\subsection{Part (a): jitter}

Let $t_j=g_j/Q+\eta_j$ with $g_j$ grid points and $\eta_j$ i.i.d.\ mean-zero with
characteristic function $\chi_\eta$, and $\nu-\omega_k=rQ$, $r\in\Z\setminus\{0\}$.  The
fold coherence is
\[
C=\frac1N\sum_je^{i2\pi(\nu-\omega_k)t_j}
 =\frac1N\sum_j\underbrace{e^{i2\pi rg_j}}_{=1}\,e^{i2\pi rQ\eta_j}
 =\frac1N\sum_je^{i2\pi rQ\eta_j},
\]
an empirical characteristic function with $\mathbb E C=\chi_\eta(2\pi rQ)$ and
concentration \eqref{eq:hoeffding}: $|C-\chi_\eta(2\pi rQ)|\le
2\sqrt{\log(4/\delta)/N}$ with probability $\ge1-\delta$.  For Gaussian jitter,
$\chi_\eta(u)=e^{-u^2\sigma_t^2/2}$, so $\chi_\eta(2\pi rQ)=e^{-2(\pi rQ\sigma_t)^2}$;
for uniform jitter of width $w$, $\chi_\eta(2\pi rQ)=\mathrm{sinc}(rQw)$.

\emph{From coherence to visibility.}  Let $\bm b=\Phimat^{*}\bm\phi_\nu/N\in\C^m$ and
$\bm G_N=\Phimat^{*}\Phimat/N$.  Then
\[
\vt(\nu)^2=\frac1N\norm{\bm\phi_\nu}^2-\frac1N\bm\phi_\nu^{*}\bm P_\Lam\bm\phi_\nu
=1-\bm b^{*}\bm G_N^{-1}\bm b .
\]
Assume additionally that the grid points $g_j$ are drawn i.i.d.\ uniformly from the grid
(a uniformly random subset drawn without replacement obeys the same bounds with the same
constants by Serfling's inequality for sampling without replacement).  Each entry of
$\bm b$ other than the fold entry, and each off-diagonal entry of $\bm G_N$, is an
empirical mean of independent unit-modulus terms whose expectation vanishes (distinct
residues modulo $Q$ average to zero over the uniform grid, and the jitter factor only
multiplies by a number of modulus $\le1$); the fold entry of $\bm b$ is $C$ above.  By
\eqref{eq:hoeffding} and a union bound over these $m+m^2$ quantities, with probability
$\ge1-\delta$ all of them are within
$\varepsilon=2\sqrt{\log(4(m+m^2)/\delta)/N}$ of their means.  Writing
$\bm G_N=\bm I+\bm E$ with $\norm{\bm E}\le m\varepsilon$ and
$\bm b=C\bm e_k+\bm u$ with $\norm{\bm u}\le\sqrt m\,\varepsilon$, a first-order
expansion of $\bm b^{*}\bm G_N^{-1}\bm b$ gives
\[
\big|\vt(\nu)^2-\big(1-|\chi_\eta(2\pi rQ)|^2\big)\big|
\ \le\ c\,m\,\varepsilon
\]
for an absolute constant $c$ (valid while $m\varepsilon<1/2$), which is the
$O_P\!\big(m\sqrt{\log(m/\delta)/N}\big)$ term of the theorem.  The small-jitter law
follows from $1-e^{-2u}\le 2u$ and $1-e^{-2u}\to 2u$ as $u\to0$ with
$u=(2\pi rQ\sigma_t)^2/2$:
$\vt(\nu)=\sqrt{1-\chi^2}+O(m\varepsilon)=2\pi|r|Q\sigma_t\,(1+o(1))$.

\subsection{Part (b): i.i.d.\ samples}

Entries of $\bm G_N-\bm G_\infty$ and of $\bm b_\nu=\Phimat^{*}\bm\phi_\nu/N$ (for every
$\nu\in\Omega$) are empirical means of independent unit-modulus variables; by
\eqref{eq:hoeffding} and a union bound over the $m^2+mK$ of them, with probability
$\ge1-\delta$ every such quantity is within
$\varepsilon=2\sqrt{\log(4(mK+m^2)/\delta)/N}$ of its limit.  Under the hypotheses, the
limiting coherences vanish, so $\norm{\bm b_\nu}\le\sqrt m\,\varepsilon$ for all
$\nu\in\Omega$, and $\norm{\bm G_N-\bm G_\infty}\le\norm{\cdot}_F\le m\varepsilon$, hence
$\lambda_{\min}(\bm G_N)\ge\lambda_{\min}-m\varepsilon>0$.  Therefore, uniformly over
$\Omega$,
\[
\at(\nu)=\norm{\bm G_N^{-1}\bm b_\nu}
\le\frac{\norm{\bm b_\nu}}{\lambda_{\min}(\bm G_N)}
\le\frac{\sqrt m\,\varepsilon}{\lambda_{\min}-m\varepsilon}. \qquad\square
\]

\begin{remark}
The bound is a union bound and is loose by design (Fig.~1b of the main paper); its role
is the \emph{rate} $N^{-1/2}$ and the uniformity over the candidate set, both matched
empirically.  It is finite exactly when $m\varepsilon<\lambda_{\min}$ and is reported as
vacuous otherwise, never silently extrapolated.
\end{remark}

\section{Proof of Theorem 4 (detection dichotomy)}

\textbf{(a)}  If $S_Tf_1=S_Tf_2$ and the observation law depends on the signal only
through $S_Tf$ (as for additive noise $\bm y=S_Tf+\bm\varepsilon$ with
$\bm\varepsilon$ independent of $f$), the distributions of $\bm y$ under the two
hypotheses are equal.  For any (possibly randomized) test $\varphi(\bm y)\in[0,1]$,
$\mathbb E_{f_1}\varphi=\mathbb E_{f_2}\varphi$: power equals size.

\textbf{(b)}  Under $\bm y=\bm D\bm\beta+\bm s+\bm\varepsilon$ with
$\bm\varepsilon\sim\mathcal N(0,\sigma^2\bm I)$ and $\bm P_D$ the projector onto the
rank-$p$ column space of $\bm D$, $(\bm I-\bm P_D)\bm y=(\bm I-\bm P_D)(\bm
s+\bm\varepsilon)$ is Gaussian with mean $(\bm I-\bm P_D)\bm s$ and covariance
$\sigma^2(\bm I-\bm P_D)$; hence $\norm{(\bm I-\bm P_D)\bm y}^2/\sigma^2$ is noncentral
$\chi^2$ with $N-p$ degrees of freedom and noncentrality
$\lambda=\norm{(\bm I-\bm P_D)\bm s}^2/\sigma^2$ (central when $\bm s=0$ or when the
component is coherent, $(\bm I-\bm P_D)\bm s=0$).  The stated power expression is the
definition of the noncentral-$\chi^2$ survival function at the central
$(1{-}\alpha)$-quantile. \hfill$\square$

\section{Background: average-case decomposition}

\begin{proposition}[average-case LS error; standard bookkeeping]
Let $\bm y=\Phimat\bm c^\star+\Phimat_{\mathrm{out}}\bm a+\bm\varepsilon$ with
$\sigma_{\min}(\Phimat)>0$, the $a_l$ zero-mean, uncorrelated, $\mathbb E|a_l|^2=p_l$,
independent of $\bm\varepsilon$, and
$\mathbb E[\bm\varepsilon\bm\varepsilon^{*}]=\sigma^2\bm I$.  Then, with
$\bm M=(\Phimat^{*}\Phimat)^{-1}\Phimat^{*}$,
\[
\mathbb E\norm{\hat{\bm c}-\bm c^\star}^2
=\sigma^2\tr\!\big((\Phimat^{*}\Phimat)^{-1}\big)
+\tr\!\big(\bm M\,\Phimat_{\mathrm{out}}\mathrm{diag}(\bm p)\Phimat_{\mathrm{out}}^{*}\,
\bm M^{*}\big).
\]
\end{proposition}
\begin{proof}
$\hat{\bm c}-\bm c^\star=\bm M(\Phimat_{\mathrm{out}}\bm a+\bm\varepsilon)$; expand the
squared norm; the cross term vanishes by independence and zero means; the two quadratic
terms are the stated traces.
\end{proof}
This proposition is retained for completeness only; it is standard omitted-variable
bias/variance bookkeeping and is not claimed as a contribution.

\end{document}
